% !TEX root =../main.tex

\subsection*{Parte A}

\begin{enumerate}[series=labsteps]

\item \textbf{Construcción de la conectividad física de la red}

\paragraph{¿Cuál es la diferencia entre un cable de UTP directo y uno cruzado?}

\subsubsection*{Evidencias}
\begin{itemize}
\item Fotografías de la topología física.
\item Fotografías de los dispositivos utilizados.
\item Evidencia del cableado realizado.
\end{itemize}


\item \textbf{Configuración básica de los equipos}

\paragraph{¿Por qué se realiza toda esta configuración básica a los dispositivos 
    Switch y Router?}

\subsubsection*{Evidencias}
\begin{itemize}
\item Capturas de la configuración en la CLI del Switch.
\item Capturas de la configuración en la CLI del Router.
\item Fotografías de la conexión por consola.
\end{itemize}


\item \textbf{Configuración del TCP/IP}

\paragraph{¿Por qué la configuración del direccionamiento IP es distinta entre un PC, 
un switch, un router y un hub?}

\paragraph{¿Por qué se asigna el direccionamiento del switch en la interfaz VLAN 1, 
mientras que al router se le asigna en otra interfaz?}

\subsubsection*{Evidencias}
\begin{itemize}
\item Capturas de la configuración de red del PC.
\item Capturas de la configuración de la interfaz VLAN del Switch.
\item Capturas de la configuración de la interfaz del Router.
\item Evidencia de la configuración del Hub.
\end{itemize}


\item \textbf{Pruebas de conectividad End-to-End (E2E)}

\paragraph{¿Se pudo realizar la prueba?}

\paragraph{¿Cómo se realizó la prueba? (Modo del dispositivo y comandos)}

\paragraph{¿Qué mensaje apareció después de hacer la prueba y qué significa?}

\paragraph{¿La prueba fue exitosa o no?}

\paragraph{¿Cuál es el propósito y utilidad de realizar estas pruebas de 
conectividad end-to-end en una red y de ICMP?}

\subsubsection*{Evidencias}
\begin{itemize}
\item Capturas de los comandos \code{ping}.
\item Resultados de las pruebas desde los distintos dispositivos.
\item Evidencia de las respuestas ICMP.
\end{itemize}

\end{enumerate}
